\section{Discussion, Ethics and Conclusion}
\label{sec:discussion}

The design gives each party what it does well: the model writes a fresh problem for this child, in the child's language and a setting the child likes; the service decides what to practise and checks the key and the repeats. The self-check is the model's opinion of its own work, and the solver its own formalisation, which the service runs and judges. Every defect that breaks a rule a check states was caught, while a question turned round under its unchanged solver, as a misreading the solver shares would leave it, and the same answer in words went through, and copies with new numbers were not reliably caught. The learner model needs the most change: it estimates a child who stays put as well as the best baseline, but its mastery rule declares mastery the child does not have, and its shrinking step leaves a child who learns or jumps behind, where Glicko-2 keeps up and a floor helps only in part; a step that does not shrink away and a stricter rule are future work.

\paragraph{Ethics.} No child took part in this study. The service learns no name, birth date or school: it knows a child by a pseudonym and a random identifier, and a parent in its logs by a keyed hash of their account, beside the IP addresses the platform logs; as a pseudonym does not make data anonymous~\cite{gdpr2016regulation}, the design keeps as little as it can: the profile lives in the parent's own Google Drive and the service keeps no copy, task texts, answers and the pseudonym stay out of logs and metrics, the service trains no model, and its policy promises no advertising. The child's grade and interests and the parent's notes do reach the chat's model with every task. Whether COPPA or Article~8 of the GDPR applies is a legal question we leave open~\cite{coppa2025rule,gdpr2016regulation}. Neither host admits a child of primary age as a user: Claude's users must be at least \stat{literature}{anthropic_minimum_age}, ChatGPT's at least \stat{literature}{openai_minimum_age}~\cite{anthropic2025consumerterms,openai2026terms}. The design assumes an adult beside the child who runs the lesson, as the service's terms say; on the card the child answers, and what the child types into the chat goes to a model whose replies no check covers. A study with children would need an ethics board's review, the parent's permission and the child's assent.

\paragraph{Cost and openness.} The service calls no model, so every task is written within the family's own use of the chat and the operator pays for computation alone, and it speaks only the open protocol, tied to no host's model. \ifanonymous The reference tasks with their solvers, the code the experiments run, their data and seeds, and the evidence ledger are in an anonymised artifact under the MIT license, \ArtifactURL.\else The code and the reference tasks are open under the MIT license at \ArtifactURL, and the experiments' code, data and seeds and the evidence ledger are archived at a stub archive.\fi

\paragraph{Conclusion.} Letting a chat model write for a child is safe only as far as what it writes is checked. We built a service that admits a model's task only after a program the model wrote has agreed with its key under both labellings and keeps the answer sealed in the child's profile until the child answers, and we measured what its checks catch and what they cannot. Acceptance in real hosts, generation by other model families, a stricter mastery rule and a study with children come next.
